% ========== CHAPTER 19 COVER: FRONTEND IMPLEMENTATION ==========
% Creative and elegant chapter introduction with sophisticated visual elements
% Following the established 4-element pattern
% Requirements: Academic research focus on frontend architecture

\vspace{2cm}

% 1. BIG CHAPTER NUMBER in content-reflective shape (component tree for React)
\begin{center}
\begin{tikzpicture}
    % React component tree structure
    % Root component (top)
    \fill[brandPrimary!15] (0,2) rectangle (0.8,2.8);
    \draw[brandPrimary!40, line width=2pt] (0,2) rectangle (0.8,2.8);
    \node[brandPrimary, font=\tiny\bfseries] at (0.4,2.4) {APP};
    
    % Second level components
    \foreach \x/\label in {-2/UI, -0.6/STATE, 0.6/COMP, 2/PERF} {
        \fill[brandSecondary!15] (\x,0.8) rectangle (\x+0.6,1.4);
        \draw[brandSecondary!40, line width=1.5pt] (\x,0.8) rectangle (\x+0.6,1.4);
        \node[brandSecondary, font=\tiny\bfseries] at (\x+0.3,1.1) {\label};
        % Connection lines from root
        \draw[brandAccent!50, line width=1pt] (0.4,2) -- (\x+0.3,1.4);
    }
    
    % Third level components (leaf nodes)
    \foreach \x in {-2.5, -1.5, -1, -0.2, 0.2, 1, 1.5, 2.5} {
        \fill[brandTertiary!15] (\x,-0.5) rectangle (\x+0.4,0.1);
        \draw[brandTertiary!40, line width=1pt] (\x,-0.5) rectangle (\x+0.4,0.1);
        % Connection lines to parent components
        \pgfmathsetmacro{\parent}{(\x < -0.8) ? -1.7 : ((\x < 0.4) ? -0.3 : ((\x < 1.2) ? 0.9 : 2.3))}
        \draw[brandAccent!30, line width=0.8pt] (\parent,0.8) -- (\x+0.2,0.1);
    }
    
    % Virtual DOM representation (dashed overlay)
    \draw[brandPrimary!20, dashed, line width=1.5pt] (-3,-1) rectangle (3,3.2);
    \node[brandPrimary!40, font=\small] at (-2.5,3) {Virtual DOM};
    
    % Chapter number (centered)
    \node[font=\fontsize{80}{80}\selectfont\bfseries, color=brandPrimary] at (0,0.8) {19};
\end{tikzpicture}
\end{center}

\vspace{1cm}

% 2. CHAPTER TOPIC with elegant typography
\begin{center}
{\fontsize{28}{32}\selectfont\textcolor{brandSecondary}{\textit{Frontend}}}\\
\vspace{0.3cm}
{\fontsize{24}{28}\selectfont\textcolor{brandTertiary}{\textit{Implementation}}}
\end{center}

\vspace{0.5cm}

% 3. NICE SEPARATOR with decorative elements
\begin{center}
\textcolor{brandAccent}{
\rule{2cm}{0.8pt} \quad
\raisebox{-0.2ex}{\Large$\square$} \quad
\rule{4cm}{0.8pt} \quad
\raisebox{-0.2ex}{\Large$\square$} \quad
\rule{2cm}{0.8pt}
}
\end{center}

\vspace{2cm}

% 4. "THIS CHAPTER EXPLORES" content box
\begin{tcolorbox}[
    enhanced,
    colback=white,
    colframe=brandPrimary!40,
    boxrule=1pt,
    arc=12pt,
    drop shadow={brandPrimary!20},
    left=1.5cm,
    right=1.5cm,
    top=1.2cm,
    bottom=1.2cm,
    before upper={\parindent0pt}
]
\begin{center}
{\Large\textbf{\textcolor{brandPrimary}{This chapter explores}}}
\end{center}

\vspace{0.8cm}

\begin{compactlist}
    \item {\large React architecture and component design for AI-first development environments}
    \item {\large Core UI component implementation and performance optimization}
    \item {\large Styling system and theming architecture for professional IDEs}
    \item {\large Tauri integration and native API communication patterns}
    \item {\large Accessibility considerations and responsive design principles}
    \item {\large Research contributions to AI-assisted development interface design}
\end{compactlist}
\end{tcolorbox}

\clearpage
